% Generated draft: requires mathematical and source review.
\documentclass[11pt]{article}
\usepackage[margin=25mm]{geometry}
\usepackage{fontspec}
\setmainfont{Latin Modern Roman}
\usepackage{amsmath,amssymb,mathtools,unicode-math}
\setmathfont{Latin Modern Math}
\usepackage{enumitem,needspace,xurl,hyperref}
\setlength{\parindent}{0pt}
\setlength{\parskip}{5pt}
\setlength{\emergencystretch}{3em}
\newcommand{\N}{\mathbb{N}}
\newcommand{\Z}{\mathbb{Z}}
\newcommand{\Q}{\mathbb{Q}}
\newcommand{\R}{\mathbb{R}}
\newcommand{\C}{\mathbb{C}}
\newcommand{\F}{\mathbb{F}}
\newcommand{\A}{\mathbb{A}}
\newcommand{\PP}{\mathbb{P}}
\newcommand{\E}{\mathbb{E}}
\newcommand{\eps}{\varepsilon}
\newcommand{\dd}{\,\mathrm{d}}
\begin{document}
\section*{UnsolvedMath ID: 30003045}
\textbf{Source problem code:} catalogue entry
\subsection{Definitions and mathematical statement}
The supplied background places this question in characteristic \(2\). Let \(F\) be a field of characteristic \(2\), let \(L/F\) be a finite purely inseparable extension, and let \(K/F\) be a separable extension of degree \(2\). Pure inseparability means that for every \(x\in L\), some power \(x^{2^e}\), with \(e\geq 0\), lies in \(F\).

Under the usual positive-fold convention, a bilinear Pfister form is a symmetric bilinear form
\[
 B=\bigotimes_{i=1}^{r}\langle 1,a_i\rangle_b,
 \qquad r\geq 1,\quad a_i\in F^\times,
\]
where \(\langle 1,a_i\rangle_b\) has diagonal Gram matrix \(\operatorname{diag}(1,a_i)\). For an extension \(E/F\), write \(B_E\) for scalar extension. Require both \(B\) and \(B_L\) to be anisotropic: \(B_E(v,v)\neq 0\) for every nonzero vector \(v\), for \(E=F,L\). The record does not specify the number \(r\) of Pfister factors or its relationship, if any, to the cohomological degree parameter.

Here is a conditional formulation using the standard characteristic-\(2\) Kato-cohomology interpretation of the supplied notation. For \(n\geq 0\), let \(\Omega_E^n\) be the absolute differential \(n\)-forms of a characteristic-\(2\) field \(E\), with \(\Omega_E^0=E\) and \(\Omega_E^{-1}=0\). Define
\[
 H_2^{n+1}(E)=\operatorname{coker}\!\left(
 \wp:\Omega_E^n\longrightarrow
 \Omega_E^n/\dd\Omega_E^{n-1}\right),
\]
where the additive map \(\wp\) is characterized on logarithmic forms by
\[
 \wp\!\left(x\,\frac{\dd b_1}{b_1}\wedge\cdots\wedge\frac{\dd b_n}{b_n}\right)
 = (x^2-x)\,\frac{\dd b_1}{b_1}\wedge\cdots\wedge\frac{\dd b_n}{b_n}
 \pmod{\dd\Omega_E^{n-1}},
 \qquad x\in E,\ b_i\in E^\times.
\]
For \(n=0\), this means \(H_2^1(E)=E/\{x^2-x:x\in E\}\). For any field extension \(E/F\), put
\[
 H_2^{n+1}(E/F)=\ker\!\left(H_2^{n+1}(F)\longrightarrow H_2^{n+1}(E)\right),
\]
using the natural scalar-extension map.

To interpret the function-field notation, provisionally let \(F(B)\) be the function field of the projective quadric defined by \(q_B(v)=B(v,v)=0\). Choose compatible embeddings into a common overfield, and interpret \(K(B)=K\mathbin{\cdot}F(B)\) and \(L\mathbin{\cdot}K(B)=L\mathbin{\cdot}K\mathbin{\cdot}F(B)\). These function-field and compositum conventions are not defined in the supplied record and require confirmation against the cited report.

For each intended admissible value of \(n\), the supplied corrected question asks whether, for every such choice of \(F,L,K,B\),
\[
 H_2^{n+1}\!\left(L\mathbin{\cdot}K(B)/F\right)
 =H_2^{n+1}(L/F)
  +H_2^{n+1}(K/F)
  +H_2^{n+1}(F(B)/F).
\]
The right-hand side is the sum of three subgroups of \(H_2^{n+1}(F)\). The source record does not state the range or quantification of \(n\); the convention \(n\geq0\) above defines the candidate notation without establishing the intended scope of the question. Its original transcription instead uses \(H^{2n+1}\), whereas its corrected statement uses \(H_2^{n+1}\). This draft follows the corrected display conditionally and does not resolve that discrepancy independently.
\subsection{Short English statement}
In characteristic \(2\), ask whether the cohomology classes killed by adjoining a finite purely inseparable extension, a separable quadratic extension, and a bilinear Pfister function field are exactly the sum of the three individual kernels, assuming the Pfister form remains anisotropic over the purely inseparable extension. The degree range and notation require confirmation; the supplied open status is dated curation.
\subsection{Sources}
\begin{itemize}
\item Supplied local UnsolvedMath registry record, permanent ID \(30003045\), title \emph{Cohomology Kernels for Inseparable Pfister Composita}, source problem code \texttt{catalogue entry}, in the Oberwolfach Reports Open Problems collection. The proposed year is recorded as 2016; no individual proposer is supplied.
\item \emph{Algebraic Cobordism and Projective Homogeneous Varieties} (2016), Oberwolfach Reports, pp.\ 215--265, DOI \texttt{10.4171/owr/2016/5}. Supplied source citation and literature-source label: original OWR report. \url{https://doi.org/10.4171/owr/2016/5}. The supplied background describes these reports as records of mathematical workshops and their open-problem discussions.
\item Supplied literature-source label: Aravire--Laghribi--O'Ryan kernel computations. \url{https://doi.org/10.1016/j.jalgebra.2015.11.022}. No article title or precise supporting passage is supplied.
\item Literature assessment dated exactly \texttt{2026-08-21}: the supplied curation labels the problem open and states that the compositum-of-kernels equality in characteristic \(2\) remains posed in the report, while published bilinear-Pfister kernel results cover important cases without establishing the full asserted equality. This is supplied dated literature triage, not independently verified here; its time-sensitive status and cited claims require independent verification.
\item The registry labels the statement \texttt{corrected\_verified} and reports that the official passage was checked and rewritten as a self-contained problem preserving its mathematical content. That verification is a supplied registry assertion; no source visit or independent mathematical verification is claimed in this draft.
\end{itemize}
\end{document}
